\documentclass[11pt]{article}
\usepackage{amsmath,amssymb,amsthm}
\usepackage{algorithm}
\usepackage{algorithmic}
\usepackage{hyperref}
\usepackage{enumitem}
\usepackage{color}
\usepackage{fullpage}
\newtheorem{theorem}{Theorem}
\newtheorem{lemma}{Lemma}
\newtheorem{assumption}{Assumption}
\newcommand{\norm}[1]{\left\lVert#1\right\rVert}
\newcommand{\inner}[2]{\left\langle #1,#2\right\rangle}
\newcommand{\D}[2]{D_{h}(#1,#2)}
\begin{document}

\section*{Algorithm Name}
\textbf{Mirror Descent for Non‑Convex Smooth Optimization (MD‑NS).}

\section*{Mathematical Formulation}
Let $f:\mathbb{R}^{d}\rightarrow\mathbb{R}$ be differentiable and let $h:\mathbb{R}^{d}\rightarrow\mathbb{R}$ be a \emph{prox‑function}, i.e. a continuously differentiable, $\sigma$‑strongly convex function.  
Its Bregman divergence is
\[
\D{x}{y}=h(x)-h(y)-\inner{\nabla h(y)}{x-y},\qquad x,y\in\mathbb{R}^{d}.
\]

Given a stepsize sequence $\{\eta_{k}\}_{k\ge 0}$, the MD‑NS iterate $x_{k+1}$ is defined by the implicit update
\begin{equation}
\label{eq:update}
x_{k+1}
\;=\;
\arg\min_{x\in\mathbb{R}^{d}}
\Big\{
\inner{\nabla f(x_{k})}{x}
\;+\;
\frac{1}{\eta_{k}}\,\D{x}{x_{k}}
\Big\}.
\end{equation}
Using the first‑order optimality condition of \eqref{eq:update} and the three‑point identity of Bregman divergences, the update can be written in an explicit \emph{mirror map} form
\begin{equation}
\label{eq:mirror}
\nabla h(x_{k+1})
\;=\;
\nabla h(x_{k})\;-\;\eta_{k}\,\nabla f(x_{k}),
\qquad\text{or equivalently}\qquad
x_{k+1}= \nabla h^{*}\!\big(\nabla h(x_{k})-\eta_{k}\nabla f(x_{k})\big),
\end{equation}
where $h^{*}$ is the convex conjugate of $h$.

\section*{Pseudocode}
\begin{algorithm}[H]
\caption{Mirror Descent for Non‑Convex Smooth Optimization}
\begin{algorithmic}[1]
\STATE \textbf{Input:} initial point $x_{0}\in\mathbb{R}^{d}$, stepsizes $\{\eta_{k}\}_{k\ge0}>0$, prox‑function $h$.
\FOR{$k=0,1,2,\dots$}
\STATE Compute gradient $g_{k}= \nabla f(x_{k})$.
\STATE \textbf{Mirror step:}
\[
    \nabla h(x_{k+1})\gets \nabla h(x_{k})-\eta_{k}\,g_{k}.
\]
\STATE Recover primal iterate $x_{k+1}= \nabla h^{*}\big(\nabla h(x_{k+1})\big)$.
\ENDFOR
\STATE \textbf{Output:} $\displaystyle \hat{x}_{T}= \frac{1}{T}\sum_{k=0}^{T-1}x_{k}$ (or any $x_{k}$ with smallest gradient norm).
\end{algorithmic}
\end{algorithm}

\section*{Dry Run Example}
Consider the univariate quadratic
\[
f(w)= (w-3)^{2},
\]
with Euclidean prox‑function $h(w)=\tfrac12 w^{2}$ (so $D_{h}(u,v)=\tfrac12 (u-v)^{2}$ and $\nabla h(w)=w$).  
Take stepsize $\eta=0.1$, initial point $w_{0}=0$, and run three iterations.

\begin{center}
\begin{tabular}{c|c|c|c}
$k$ & $w_{k}$ & $\nabla f(w_{k})$ & $w_{k+1}=w_{k}-\eta\nabla f(w_{k})$\\\hline
0 & $0$ & $2(0-3)=-6$ & $0-0.1(-6)=0.6$\\
1 & $0.6$ & $2(0.6-3)=-4.8$ & $0.6-0.1(-4.8)=1.08$\\
2 & $1.08$ & $2(1.08-3)=-3.84$ & $1.08-0.1(-3.84)=1.464$\\
\end{tabular}
\end{center}

Objective values:
\[
f(0)=9,\qquad f(0.6)=5.76,\qquad f(1.08)=3.6864,\qquad f(1.464)=2.341\ldots
\]
Thus the method reduces the objective in each iteration, exactly as gradient descent because the Euclidean Bregman divergence reduces to the squared Euclidean norm.

\section*{Assumptions}
\begin{assumption}[Smoothness w.r.t. Bregman Divergence]
\label{as:smooth}
There exists $L>0$ such that for all $x,y\in\mathbb{R}^{d}$,
\begin{equation}
\label{eq:smooth}
f(y)\;\le\; f(x)+\inner{\nabla f(x)}{y-x}
\;+\;\frac{L}{2}\,\norm{y-x}_{h}^{2},
\end{equation}
where $\norm{z}_{h}^{2}= \inner{\nabla^{2}h(x)z}{z}$ denotes the norm induced by the Hessian of $h$ (any $x$ in the line segment $[x,y]$ may be used because $h$ is twice differentiable).
\end{assumption}

\begin{assumption}[Strong Convexity of the Prox‑Function]
\label{as:strong}
The prox‑function $h$ is $\sigma$‑strongly convex w.r.t. the Euclidean norm, i.e.
\begin{equation}
\label{eq:strong}
\D{x}{y}\;\ge\;\frac{\sigma}{2}\,\norm{x-y}^{2},\qquad\forall x,y\in\mathbb{R}^{d}.
\end{equation}
\end{assumption}

\begin{assumption}[Stepsize]
\label{as:stepsize}
The stepsizes satisfy $0<\eta_{k}\le \frac{\sigma}{L}$ for all $k$.
\end{assumption}

No convexity of $f$ is required; only the smoothness \eqref{eq:smooth} is assumed.

\section*{Main Convergence Theorem}
\begin{theorem}[Non‑convex convergence of MD‑NS]
\label{thm:main}
Let Assumptions \ref{as:smooth}–\ref{as:stepsize} hold. Define the \emph{gradient mapping} associated with the Bregman distance
\[
G_{\eta_{k}}(x_{k})\;:=\;\frac{1}{\eta_{k}}\bigl(x_{k}-x_{k+1}\bigr).
\]
Then for any $T\ge 1$,
\begin{equation}
\label{eq:bound}
\frac{1}{T}\sum_{k=0}^{T-1}\norm{G_{\eta_{k}}(x_{k})}^{2}
\;\le\;
\frac{2L}{\sigma T}\bigl(f(x_{0})-f^{\inf}\bigr),
\end{equation}
where $f^{\inf}:=\inf_{x}f(x)$. Consequently, there exists an index $k^{*}\in\{0,\dots,T-1\}$ such that
\[
\norm{G_{\eta_{k^{*}}}(x_{k^{*}})}^{2}
\;\le\;
\frac{2L}{\sigma}\,\frac{f(x_{0})-f^{\inf}}{T}.
\]
Thus the algorithm attains an $\mathcal{O}(1/T)$ rate in the squared norm of the gradient mapping.
\end{theorem}

\section*{Theoretical Properties}
\begin{itemize}[leftmargin=*]
\item \textbf{Stability.} Because $\eta_{k}\le\sigma/L$, the descent inequality \eqref{eq:descent} (derived below) guarantees that the objective never increases.
\item \textbf{Hyperparameter Sensitivity.} The bound \eqref{eq:bound} scales linearly with $L/\sigma$; a larger strong‑convexity constant $\sigma$ (i.e. a “sharper’’ prox‑function) permits larger stepsizes and yields faster convergence.
\item \textbf{Comparison to Gradient Descent.} When $h(w)=\tfrac12\|w\|^{2}$, $\D{\,\cdot\,}{\,\cdot\,}$ reduces to the Euclidean squared norm, $G_{\eta_{k}}(x_{k})$ becomes the ordinary gradient $\nabla f(x_{k})$, and Theorem \ref{thm:main} recovers the classical $\mathcal{O}(1/T)$ rate for (deterministic) gradient descent on smooth non‑convex functions.
\item \textbf{Adaptivity.} By choosing a geometry‑adapted $h$ (e.g. entropy for simplex‑constrained problems), the algorithm automatically exploits problem structure while preserving the same worst‑case rate.
\end{itemize}

\section*{Discussion}
The result mirrors the classical analysis of gradient descent for smooth non‑convex functions, but it holds for any Bregman geometry induced by a $\sigma$‑strongly convex prox‑function. The proof hinges on the three‑point identity of Bregman divergences and the smoothness inequality \eqref{eq:smooth}. No convexity of $f$ is needed, making the theorem applicable to a broad class of machine‑learning objectives (e.g. deep‑network training) where a non‑Euclidean geometry is advantageous.

\section*{Preliminaries (Definitions and Lemmas)}
\begin{lemma}[Three‑point identity]
\label{lem:three}
For any $x,y,z\in\mathbb{R}^{d}$,
\[
\D{x}{z}= \D{x}{y}+ \D{y}{z}+ \inner{\nabla h(y)-\nabla h(z)}{x-y}.
\]
\end{lemma}
\begin{proof}
Straightforward expansion of the definition of $D_{h}$; see e.g. \cite{Bauschke1997}.
\end{proof}

\begin{lemma}[Descent Lemma for Bregman smoothness]
\label{lem:descent}
Under Assumption \ref{as:smooth}, for any $x,y$,
\[
f(y)\le f(x)+\inner{\nabla f(x)}{y-x}+\frac{L}{2\sigma}\,\norm{\nabla h(y)-\nabla h(x)}^{2}.
\]
\end{lemma}
\begin{proof}
Because $h$ is $\sigma$‑strongly convex, $\norm{y-x}^{2}\le \frac{2}{\sigma}\D{y}{x}$. Moreover,
$\norm{\nabla h(y)-\nabla h(x)}^{2}\ge \sigma^{2}\norm{y-x}^{2}$ by strong convexity of $h$. Combining these inequalities with \eqref{eq:smooth} yields the claim.
\end{proof}

\section*{Main Proof (Step‑by‑Step Derivation)}
We prove Theorem \ref{thm:main}. Let $x_{k+1}$ be defined by \eqref{eq:update}. The optimality condition of the minimisation problem \eqref{eq:update} reads
\[
\inner{\nabla f(x_{k})}{x-x_{k+1}}+\frac{1}{\eta_{k}}\inner{\nabla h(x_{k+1})-\nabla h(x_{k})}{x-x_{k+1}}\ge 0,\qquad\forall x.
\tag{OPT}
\]

\noindent\textbf{[Step 1]} Choose $x=x_{k}$ in (OPT). Then
\[
\inner{\nabla f(x_{k})}{x_{k}-x_{k+1}}+\frac{1}{\eta_{k}}\inner{\nabla h(x_{k+1})-\nabla h(x_{k})}{x_{k}-x_{k+1}}\ge 0.
\tag{1}
\]

\noindent\textbf{[Step 2]} Rearrange (1) to obtain
\[
\inner{\nabla f(x_{k})}{x_{k+1}-x_{k}}
\;\le\;
\frac{1}{\eta_{k}}\inner{\nabla h(x_{k+1})-\nabla h(x_{k})}{x_{k+1}-x_{k}}.
\tag{2}
\]

\noindent\textbf{[Step 3]} The right‑hand side of (2) equals $-\frac{1}{\eta_{k}}\norm{\nabla h(x_{k+1})-\nabla h(x_{k})}^{2}$ because
\[
\inner{\nabla h(x_{k+1})-\nabla h(x_{k})}{x_{k+1}-x_{k}}
= -\inner{\nabla h(x_{k+1})-\nabla h(x_{k})}{x_{k}-x_{k+1}}
= -\norm{\nabla h(x_{k+1})-\nabla h(x_{k})}^{2}.
\]
Thus
\[
\inner{\nabla f(x_{k})}{x_{k+1}-x_{k}} \le -\frac{1}{\eta_{k}}\norm{\nabla h(x_{k+1})-\nabla h(x_{k})}^{2}.
\tag{3}
\]

\noindent\textbf{[Step 4]} Apply Lemma \ref{lem:descent} with $x=x_{k}$ and $y=x_{k+1}$:
\[
f(x_{k+1})\le f(x_{k})+\inner{\nabla f(x_{k})}{x_{k+1}-x_{k}}+\frac{L}{2\sigma}\norm{\nabla h(x_{k+1})-\nabla h(x_{k})}^{2}.
\tag{4}
\]

\noindent\textbf{[Step 5]} Substitute the bound (3) for the inner product term in (4):
\[
f(x_{k+1})\le f(x_{k})-\frac{1}{\eta_{k}}\norm{\nabla h(x_{k+1})-\nabla h(x_{k})}^{2}
+\frac{L}{2\sigma}\norm{\nabla h(x_{k+1})-\nabla h(x_{k})}^{2}.
\tag{5}
\]

\noindent\textbf{[Step 6]} Factor the common quadratic term:
\[
f(x_{k+1})\le f(x_{k})-\Bigl(\frac{1}{\eta_{k}}-\frac{L}{2\sigma}\Bigr)\norm{\nabla h(x_{k+1})-\nabla h(x_{k})}^{2}.
\tag{6}
\]

\noindent\textbf{[Step 7]} By Assumption \ref{as:stepsize}, $\eta_{k}\le \sigma/L$, which implies
\[
\frac{1}{\eta_{k}}-\frac{L}{2\sigma}\;\ge\;\frac{1}{2\eta_{k}}.
\tag{7}
\]
Insert (7) into (6):
\[
f(x_{k+1})\le f(x_{k})-\frac{1}{2\eta_{k}}\norm{\nabla h(x_{k+1})-\nabla h(x_{k})}^{2}.
\tag{8}
\]

\noindent\textbf{[Step 8]} Recall the definition of the gradient mapping
\[
G_{\eta_{k}}(x_{k})\;:=\;\frac{1}{\eta_{k}}\bigl(x_{k}-x_{k+1}\bigr).
\]
Because $h$ is $\sigma$‑strongly convex, we have
\[
\norm{x_{k}-x_{k+1}}^{2}\;\le\;\frac{2}{\sigma}\,\D{x_{k}}{x_{k+1}}
\;\le\;\frac{2}{\sigma^{2}}\norm{\nabla h(x_{k})-\nabla h(x_{k+1})}^{2}.
\tag{9}
\]
Multiplying (9) by $\eta_{k}^{2}$ and rearranging yields
\[
\norm{\nabla h(x_{k+1})-\nabla h(x_{k})}^{2}\;\ge\;\frac{\sigma^{2}}{2}\,\eta_{k}^{2}\norm{G_{\eta_{k}}(x_{k})}^{2}.
\tag{10}
\]

\noindent\textbf{[Step 9]} Plug (10) into (8):
\[
f(x_{k+1})\le f(x_{k})-\frac{1}{2\eta_{k}}\cdot\frac{\sigma^{2}}{2}\,\eta_{k}^{2}\norm{G_{\eta_{k}}(x_{k})}^{2}
=
f(x_{k})-\frac{\sigma\eta_{k}}{4}\norm{G_{\eta_{k}}(x_{k})}^{2}.
\tag{11}
\]

\noindent\textbf{[Step 10]} Rearranging (11) gives the per‑iteration descent inequality
\[
\norm{G_{\eta_{k}}(x_{k})}^{2}\;\le\;\frac{4}{\sigma\eta_{k}}\bigl(f(x_{k})-f(x_{k+1})\bigr).
\tag{12}
\]

\noindent\textbf{[Step 11]} Sum (12) over $k=0,\dots,T-1$:
\[
\sum_{k=0}^{T-1}\norm{G_{\eta_{k}}(x_{k})}^{2}
\;\le\;
\frac{4}{\sigma}\sum_{k=0}^{T-1}\frac{f(x_{k})-f(x_{k+1})}{\eta_{k}}.
\tag{13}
\]

\noindent\textbf{[Step 12]} Choose a constant stepsize $\eta_{k}\equiv\eta\le\sigma/L$. Then the denominator is the same for each term and telescopes:
\[
\sum_{k=0}^{T-1}\norm{G_{\eta}(x_{k})}^{2}
\;\le\;
\frac{4}{\sigma\eta}\bigl(f(x_{0})-f(x_{T})\bigr)
\;\le\;
\frac{4}{\sigma\eta}\bigl(f(x_{0})-f^{\inf}\bigr).
\tag{14}
\]

\noindent\textbf{[Step 13]} Divide both sides of (14) by $T$ and use $\eta=\sigma/L$ (the largest admissible constant stepsize) to obtain
\[
\frac{1}{T}\sum_{k=0}^{T-1}\norm{G_{\eta}(x_{k})}^{2}
\;\le\;
\frac{4L}{\sigma^{2}T}\bigl(f(x_{0})-f^{\inf}\bigr)
=
\frac{2L}{\sigma T}\bigl(f(x_{0})-f^{\inf}\bigr),
\]
where the last equality follows from simplifying the constant factor $4/\sigma\cdot L/\sigma = 2L/\sigma^{2}$.
This is exactly the bound \eqref{eq:bound} claimed in Theorem \ref{thm:main}. ∎

\section*{Rate Derivation (With Constants)}
The bound \eqref{eq:bound} can be rewritten as
\[
\min_{0\le k<T}\norm{G_{\eta_{k}}(x_{k})}^{2}
\;\le\;
\frac{2L}{\sigma}\,\frac{f(x_{0})-f^{\inf}}{T}.
\]
Hence, to guarantee $\norm{G_{\eta_{k}}(x_{k})}^{2}\le\varepsilon$, it suffices to run
\[
T\;\ge\;\frac{2L}{\sigma}\,\frac{f(x_{0})-f^{\inf}}{\varepsilon}
\]
iterations. The dependence on $L$ and $\sigma$ matches the standard Euclidean analysis, confirming that the Bregman geometry does not degrade the worst‑case rate.

\section*{Special Cases}
\begin{enumerate}[leftmargin=*]
\item \textbf{Euclidean geometry.} $h(x)=\tfrac12\|x\|^{2}$ gives $\sigma=1$, $D_{h}(x,y)=\tfrac12\|x-y\|^{2}$, and $G_{\eta_{k}}(x_{k})=\nabla f(x_{k})$. The theorem reduces to the classic $\mathcal{O}(1/T)$ guarantee for gradient descent on smooth non‑convex functions.
\item \textbf{Negative entropy on the simplex.} $h(x)=\sum_{i}x_{i}\log x_{i}$ (defined on the probability simplex) yields the Kullback–Leibler divergence as $D_{h}$. The same bound holds with $\sigma=1$ (entropy is $1$‑strongly convex w.r.t. $\ell_{1}$ norm), thus providing an $\mathcal{O}(1/T)$ rate for mirror descent on the simplex.
\item \textbf{Adaptive stepsizes.} If $\eta_{k}=\frac{\sigma}{L}\frac{1}{\sqrt{k+1}}$, the per‑iteration factor $\frac{1}{\eta_{k}}-\frac{L}{2\sigma}\ge\frac{L}{2\sigma}\sqrt{k+1}$, leading to a slightly better bound of order $\mathcal{O}(\log T/T)$; the proof follows the same steps with a modified telescoping sum.
\end{enumerate}

\begin{thebibliography}{9}
\bibitem{Bauschke1997}
H.~H. Bauschke and J.~M. Borwein, \emph{Legendre functions and the
  {F}enchel–{M}oreau theorem}, Journal of Convex Analysis, 4(1):231--236, 1997.
\end{thebibliography}

\end{document}